\documentclass[sigconf]{acmart}
\AtBeginDocument{%
  \providecommand\BibTeX{{%
    Bib\TeX}}}

\setcopyright{acmlicensed}
\copyrightyear{2018}
\acmYear{2018}
\acmDOI{XXXXXXX.XXXXXXX}
%% These commands are for a PROCEEDINGS abstract or paper.
\acmConference[Conference acronym 'XX]{Make sure to enter the correct
  conference title from your rights confirmation email}{June 03--05,
  2018}{Woodstock, NY}
%%
%%  Uncomment \acmBooktitle if the title of the proceedings is different
%%  from ``Proceedings of ...''!
%%
%%\acmBooktitle{Woodstock '18: ACM Symposium on Neural Gaze Detection,
%%  June 03--05, 2018, Woodstock, NY}
\acmISBN{978-1-4503-XXXX-X/2018/06}
\begin{document}

\title{Replicação AgentRx: Diagnóstico em Falhas de Sistemas Multiagentes utilizando Telemetria Estruturada}


\author{Pedro Henrique Silva Da Costa Gregório}
\email{pedro.gregorio@copin.ufcg.edu.br}
\affiliation{%
	\institution{Universidade Federal de Campina Grande (UFCG)}
	\city{Campina Grande}
	\state{PB}
	\country{Brasil}
}

\begin{abstract}
	A depuração de Sistemas Multiagentes (MAS) baseados em Grandes Modelos de Linguagem (LLMs) é um desafio crítico devido à natureza probabilística, de longo horizonte e mediada por saídas ruidosas de ferramentas dessas execuções. O estudo original do AgentRx introduziu um framework automatizado promissor que localiza a etapa crítica de uma falha e categoriza sua causa raiz, mas depende de logs textuais heterogêneos e desestruturados, exigindo normalizações complexas e sujeitas a ruídos. A premissa deste estudo de replicação é verificar se a injeção de uma entrada mais estruturada, fundamentada nos padrões de observabilidade do OpenTelemetry, consegue melhorar a classificação e acurácia diagnóstica do framework. Diferente do estudo original, que amostrou logs de execuções passadas no mundo real, o processo metodológico adotado aqui baseia-se em uma experimentação laboratorial controlada. Utilizando um benchmark sintético em um ambiente MAS simulado, o experimento avalia dois cenários distintos de representação de trajetórias (logs textuais simplificados versus telemetria estruturada), visando não apenas elevar a precisão do diagnóstico, mas também viabilizar a estimativa inédita do custo operacional desperdiçado pelas falhas.
\end{abstract}

\begin{CCSXML}
	<ccs2012>
	<concept>
	<concept_id>00000000.0000000.0000000</concept_id>
	<concept_desc>Do Not Use This Code, Generate the Correct Terms for Your Paper</concept_desc>
	<concept_significance>500</concept_significance>
	</concept>
	<concept>
	<concept_id>00000000.00000000.00000000</concept_id>
	<concept_desc>Do Not Use This Code, Generate the Correct Terms for Your Paper</concept_desc>
	<concept_significance>300</concept_significance>
	</concept>
	<concept>
	<concept_id>00000000.00000000.00000000</concept_id>
	<concept_desc>Do Not Use This Code, Generate the Correct Terms for Your Paper</concept_desc>
	<concept_significance>100</concept_significance>
	</concept>
	<concept>
	<concept_id>00000000.00000000.00000000</concept_id>
	<concept_desc>Do Not Use This Code, Generate the Correct Terms for Your Paper</concept_desc>
	<concept_significance>100</concept_significance>
	</concept>
	</ccs2012>
\end{CCSXML}

\ccsdesc[500]{Do Not Use This Code~Generate the Correct Terms for Your Paper}
\ccsdesc[300]{Do Not Use This Code~Generate the Correct Terms for Your Paper}
\ccsdesc{Do Not Use This Code~Generate the Correct Terms for Your Paper}
\ccsdesc[100]{Do Not Use This Code~Generate the Correct Terms for Your Paper}

\received{20 February 2007}
\received[revised]{12 March 2009}
\received[accepted]{5 June 2009}

\maketitle

\keywords{Sistemas Multiagentes, Telemetria, Injeção de Falhas, Agent2Agent, Model Context Protocol, Observabilidade de LLM}

\section{Introdução}
Os Sistemas Multiagentes (MAS) baseados em LLMs estão cada vez mais autônomos, executando tarefas complexas em ambientes do mundo real.
No entanto, a mesma autonomia que viabiliza essa produtividade torna esses sistemas extremamente difíceis de depurar.
Como as trajetórias são longas e a saída das ferramentas pode ser ruidosa, pequenas falhas em uma etapa inicial podem se propagar silenciosamente até que o sistema colapse de forma irrecuperável [2].

Para solucionar o problema do diagnóstico manual — que é caro e de difícil escalabilidade [3] —, o estudo de Barke et al. propôs o \textit{AgentRx},
um framework automatizado que identifica a etapa da primeira falha crítica (irrecuperável) e sua categoria raiz [4].
O estudo original foi conduzido extraindo logs desestruturados em texto puro de três domínios heterogêneos (como o $\tau$-bench e o Magentic-One), necessitando de grande esforço processual para normalizar os rastros em uma Representação Intermediária (IR) a fim de que um LLM-como-Juiz avaliasse as violações [5].

Embora o AgentRx represente o estado da arte na atribuição de falhas, a sua dependência de logs textuais ambíguos o torna suscetível a falsos positivos [6]. Este estudo de replicação propõe introduzir no processo de formação de dados um modelo de telemetria nativamente estruturado, baseado nos preceitos do OpenTelemetry [7].
O OpenTelemetry é um framework de observabilidade de código aberto que padroniza a coleta de dados de telemetria (rastreamentos, métricas e logs), desvinculando as aplicações da infraestrutura de rede e garantindo informações detalhadas de causa raiz sem ambiguidades [7, 8].
Ao invés do AgentRx precisar interpretar longas mensagens de texto para deduzir um erro, o modelo de telemetria fornecerá eventos padronizados (como \textit{timeout\_detected} ou \textit{invalid\_schema}) [9], reduzindo o esforço cognitivo do LLM-Juiz e purificando a análise de causa e efeito.

Com base nesta premissa, as principais Questões de Pesquisa (RQs) que guiam este estudo são [10]:
\begin{itemize}
	\item \textbf{RQ1:} A telemetria estruturada melhora a identificação do \textit{critical failure step} em comparação com logs textuais?
	\item \textbf{RQ2:} A identificação exata da etapa crítica permite estimar com precisão o custo operacional desperdiçado (tokens, latência e tempo) após a propagação da falha?
\end{itemize}

O restante do artigo está estruturado da seguinte forma: A Seção 2 detalha a metodologia de funcionamento do AgentRx, a arquitetura da telemetria proposta e o formato de integração de dados. A Seção 3 apresenta o planejamento e o design do experimento controlado.

\section{Metodologia}

\subsection{O Framework AgentRx}
O processo de diagnóstico do AgentRx opera recebendo como entrada o esquema das ferramentas, políticas de domínio e a trajetória falha. Como as execuções possuem formatos distintos, o AgentRx primeiramente as normaliza em uma Representação Intermediária (IR) comum. A partir disso, o sistema sintetiza "Restrições Globais" (baseadas nas políticas) e "Restrições Dinâmicas" (condicionadas ao histórico da trajetória) [11]. O framework avalia essas restrições passo a passo (utilizando validações em Python e verificações semânticas via LLM) [12]. Caso uma regra seja violada, o evento é adicionado a um log de validação auditável. Por fim, um LLM-como-Juiz analisa essas violações cruzando-as com um \textit{checklist} taxonômico de falhas para determinar o exato passo crítico e a categoria do erro [13, 14].

\subsection{Ambiente Simulado e Estruturação de Dados}
Diferente da abordagem original que dependia de benchmarks de terceiros, a metodologia desta replicação baseia-se na construção de uma Prova de Conceito (PoC) implementada com o framework LangGraph. A unidade geradora de dados é um MAS contendo quatro agentes: Coordenador, Pesquisador, Executor e Avaliador [15].

A grande alteração metodológica é a instrumentação nativa com os conceitos do \textit{OpenTelemetry}. Os logs gerados seguirão a seguinte padronização [9, 16, 17]:
\begin{itemize}
	\item \textbf{Trace:} Representa a trajetória completa da tarefa (ex: \textit{trace\_id}, \textit{start\_time}, \textit{status}).
	\item \textbf{Spans:} Segmentos detalhados que formam o Trace, como \textit{Workflow Span} (fluxo principal), \textit{Agent Span} (execução de um agente isolado) e \textit{Tool Span} (chamada de ferramenta, evidenciando \textit{tool\_input} e \textit{tool\_output}). A relação entre Spans (\textit{parent\_span\_id}) preserva a hierarquia da execução.
	\item \textbf{Events \& Attributes:} Disparam gatilhos específicos como \textit{timeout\_detected} e injetam atributos de falha reais (\textit{error.type}, \textit{error.message}).
	\item \textbf{Metrics:} Indicadores quantitativos de custos (como \textit{tokens\_in}, \textit{tokens\_out}, \textit{latency\_ms} e \textit{retry\_count}).
\end{itemize}

\subsection{Adaptação ao Fluxo do AgentRx}
Para integrar nosso modelo à arquitetura metodológica do artigo base, desenvolvemos um adaptador (\textit{Adapter}) de telemetria. Este componente é encarregado de traduzir a árvore de Traces e Spans rigorosamente para a \textit{Trajectory IR} esperada pelo pipeline original do AgentRx. Graças à formatação estruturada na origem, o gerador de restrições dinâmicas do AgentRx será reconfigurado para ler gatilhos explícitos diretamente no JSON de telemetria em vez de precisar realizar interpretações textuais extensivas [18].

\section{Análise de Experimento}

\subsection{Design de Experimentos}
O objetivo central do experimento é avaliar estatisticamente as relações de causa e efeito das representações de log no desempenho diagnóstico do framework. Baseado nos princípios de experimentação quantitativa para analisar inferências estatísticas e mitigar vieses [19], definimos um desenho de experimento \textbf{fatorial completo $2 \times 5$} [20, 21].

\textbf{Variável Independente (Fator):}
O experimento avalia dois níveis de representação da trajetória:
\begin{enumerate}
	\item \textit{Logs textuais simplificados (Baseline):} Simula a representação original, extraindo a trajetória textualmente e ocultando indicadores estruturados das métricas do OpenTelemetry.
	\item \textit{Telemetria Estruturada (Tratamento):} Dados complexos contendo os Spans, Eventos e atributos predefinidos pelo nosso modelo [10].
\end{enumerate}

Para o controle experimental, utilizamos a técnica de injeção de falhas. Foram selecionados 5 níveis de categorias da taxonomia original do AgentRx [21]: \textit{System Failure}, \textit{Invalid Invocation}, \textit{Instruction/Plan Adherence Failure}, \textit{Tool Misuse} (uso incorreto de ferramentas) e \textit{Validation Failure}.

\textbf{Variáveis Dependentes:}
As métricas utilizadas para medir o efeito dos tratamentos serão [10]:
\begin{enumerate}
	\item \textbf{Critical Step Accuracy:} Percentual de precisão na identificação do índice da etapa exata do erro.
	\item \textbf{Failure Category Accuracy:} Precisão na classificação da categoria raiz.
	\item \textbf{Average Step Distance:} Distância média, em passos, entre a predição do modelo e o erro real.
	\item \textbf{Custo Operacional Desperdiçado:} Calculado com base nas métricas de tokens e latência pós-falha a partir do marco crítico estabelecido pelo AgentRx.
\end{enumerate}

\subsection{Execução e Análise Estatística}
Para garantir relevância e balanceamento [22], cada combinação do desenho fatorial gerará 5 trajetórias simuladas independentes, resultando em 50 execuções. Para mitigar o não-determinismo inerente às avaliações de LLM, a ferramenta do AgentRx julgará cada trajetória 3 vezes [23]. O \textit{ground truth} de avaliação será extraído diretamente pelo operador de injeção de falha que gerou o respectivo problema [18].

As análises estatísticas utilizarão o cálculo de Média e Desvio Padrão para os resultados globais. Para validar a hipótese nula ($H_0$) — de que não há diferença de precisão entre os formatos de log —, aplicaremos inicialmente o Teste de normalidade de \textit{Shapiro-Wilk}. Considerando amostras dependentes executadas sob a mesma base, o contraste de hipótese será aferido via teste de \textit{Wilcoxon} pareado, adotando um nível de significância de $\alpha = 0.05$ [18].

\end{document}
\endinput

